
\documentclass{article}
\usepackage{colortbl}
\usepackage{makecell}
\usepackage{multirow}
\usepackage{supertabular}

\begin{document}

\newcounter{utterance}

\centering \large Interaction Transcript for game `function\_detective', experiment `string\_passive\_labeled\_pairs', episode 123 with Qwen3.6{-}35B{-}A3B{-}FP8{-}without{-}reasoning.
\vspace{24pt}

{ \footnotesize  \setcounter{utterance}{1}
\setlength{\tabcolsep}{0pt}
\begin{supertabular}{c@{$\;$}|p{.15\linewidth}@{}p{.15\linewidth}p{.15\linewidth}p{.15\linewidth}p{.15\linewidth}p{.15\linewidth}}
    \# & \multicolumn{2}{c}{Player} && \multicolumn{2}{c}{Game Master} \\
    \hline

    \theutterance \stepcounter{utterance}  
    & & & \multicolumn{4}{p{0.6\linewidth}}{
        \cellcolor[rgb]{0.9,0.9,0.9}{
            \makecell[{{p{\linewidth}}}]{
                \texttt{\tiny{[P1$\langle$GM]}}
                \texttt{I have a hidden Python function for you to figure out.} \\
\\ 
\texttt{Your job is to infer what it does and then submit a correct Python implementation.} \\
\\ 
\texttt{Function signature:} \\
\texttt{(text: str) {-}> str} \\
\\ 
\texttt{You have 15 turns total.} \\
\\ 
\texttt{I reveal candidate input{-}output pairs one at a time and tell you whether each one is correct.} \\
\\ 
\texttt{You must make your best SOLVE attempt no later than the second{-}to{-}last turn.} \\
\texttt{Do not wait until the very end.} \\
\\ 
\texttt{You can use these actions:} \\
\texttt{{-} JSON format described at the end of the prompt} \\
\texttt{{-} SOLVE: ```python ... ```} \\
\\ 
\texttt{Do not add import statements in your guess.} \\
\texttt{Anything you might need is already available, so just write the function body in hypothesis and return the answer.} \\
\\ 
\texttt{Preloaded examples:} \\
\texttt{1. (('   ', '   '), True)} \\
\\ 
\texttt{Scoring:} \\
\texttt{{-} correctness is required for any non{-}zero score} \\
\texttt{{-} in interactive modes, solving with fewer information requests is better} \\
\\ 
\texttt{Input formatting reminders:} \\
\texttt{{-} strings must be quoted} \\
\texttt{{-} booleans must be True/False} \\
\texttt{{-} use commas between multiple values} \\
\texttt{{-} do not write any extra action tags in the same message} \\
\\ 
\texttt{Format example:} \\
\texttt{You write:} \\
\texttt{\{'current\_mode': <probing/confirming/don't know; choose whichever best describes your current state in the task>, 'input\_rationale': <brief description of why you chose this hypothesis>, 'hypothesis': <either the exact string "unknown" if you have no guess yet OR your current best guess as runnable Python matching the given signature, with a one{-}line docstring stating what it does. Never write placeholder code, if you cannot commit to a guess, use "unknown". Match the given signature.} \\
\texttt{e.g.} \\
\texttt{def f(x):} \\
\texttt{	"""Returns the absolute value of x."""} \\
\texttt{	return x if x >= 0 else {-}x} \\
\\ 
\texttt{Use "unknown" if you have no guess.>, 'confidence\_score': <score between 1{-}100. 100 means highest confidence in your hypothesis>\}} \\
\\ 
\texttt{Output only one valid json per turn.} \\
            }
        }
    }
     \\ \\

    \theutterance \stepcounter{utterance}  
    & \multicolumn{4}{p{0.6\linewidth}}{
        \cellcolor[rgb]{0.561,0.737,0.561}{
            \makecell[{{p{\linewidth}}}]{
                \texttt{\tiny{[P1$\rangle$GM]}}
                \texttt{\{'current\_mode': 'probing', 'input\_rationale': 'The single preloaded example shows that a string of spaces compared to itself returns True. This is too vague to determine the function's logic. I need more examples to understand what transformation or comparison is being performed. I will try a few different inputs to see how the function behaves.', 'hypothesis': 'unknown', 'confidence\_score': 10\}} \\
            }
        }
    }
    & & \\ \\

    \theutterance \stepcounter{utterance}  
    & & & \multicolumn{4}{p{0.6\linewidth}}{
        \cellcolor[rgb]{0.9,0.9,0.9}{
            \makecell[{{p{\linewidth}}}]{
                \texttt{\tiny{[P1$\langle$GM]}}
                \texttt{Game Violation: ParseError: Response is not valid JSON.} \\
\texttt{Output ONLY:} \\
\texttt{One JSON with confidence\_score, current\_mode, input\_rationale, and hypothesis keys ...} \\
\texttt{OR} \\
\texttt{SOLVE: ```python ... ```} \\
            }
        }
    }
     \\ \\

    \theutterance \stepcounter{utterance}  
    & \multicolumn{4}{p{0.6\linewidth}}{
        \cellcolor[rgb]{0.561,0.737,0.561}{
            \makecell[{{p{\linewidth}}}]{
                \texttt{\tiny{[P1$\rangle$GM]}}
                \texttt{\{} \\
\texttt{"current\_mode": "probing",} \\
\texttt{"input\_rationale": "The single preloaded example shows that a string of spaces compared to itself returns True. This is too vague to determine the function's logic. I need more examples to understand what transformation or comparison is being performed. I will try a few different inputs to see how the function behaves.",} \\
\texttt{"hypothesis": "unknown",} \\
\texttt{"confidence\_score": 10} \\
\texttt{\}} \\
            }
        }
    }
    & & \\ \\

    \theutterance \stepcounter{utterance}  
    & & & \multicolumn{4}{p{0.6\linewidth}}{
        \cellcolor[rgb]{0.9,0.9,0.9}{
            \makecell[{{p{\linewidth}}}]{
                \texttt{\tiny{[P1$\langle$GM]}}
                \texttt{OUTPUT: (('a{-}b', 'n{-}o\_alt'), False)} \\
            }
        }
    }
     \\ \\

    \theutterance \stepcounter{utterance}  
    & \multicolumn{4}{p{0.6\linewidth}}{
        \cellcolor[rgb]{0.561,0.737,0.561}{
            \makecell[{{p{\linewidth}}}]{
                \texttt{\tiny{[P1$\rangle$GM]}}
                \texttt{\{} \\
\texttt{"current\_mode": "probing",} \\
\texttt{"input\_rationale": "The first example was ('   ', '   ') {-}> True. The second is ('a{-}b', 'n{-}o\_alt') {-}> False. This suggests the function might be checking if two strings are identical, or perhaps if they are anagrams, or if one is a transformation of the other. 'a{-}b' and 'n{-}o\_alt' are clearly different. Let's try a case where strings are identical but not spaces, and a case where they are anagrams.",} \\
\texttt{"hypothesis": "unknown",} \\
\texttt{"confidence\_score": 10} \\
\texttt{\}} \\
            }
        }
    }
    & & \\ \\

    \theutterance \stepcounter{utterance}  
    & & & \multicolumn{4}{p{0.6\linewidth}}{
        \cellcolor[rgb]{0.9,0.9,0.9}{
            \makecell[{{p{\linewidth}}}]{
                \texttt{\tiny{[P1$\langle$GM]}}
                \texttt{OUTPUT: (('a', 'n'), True)} \\
            }
        }
    }
     \\ \\

    \theutterance \stepcounter{utterance}  
    & \multicolumn{4}{p{0.6\linewidth}}{
        \cellcolor[rgb]{0.561,0.737,0.561}{
            \makecell[{{p{\linewidth}}}]{
                \texttt{\tiny{[P1$\rangle$GM]}}
                \texttt{\{} \\
\texttt{"current\_mode": "probing",} \\
\texttt{"input\_rationale": "We have three examples so far:\textbackslash{}n1. ('   ', '   ') {-}> True\textbackslash{}n2. ('a{-}b', 'n{-}o\_alt') {-}> False\textbackslash{}n3. ('a', 'n') {-}> True\textbackslash{}n\textbackslash{}nThe first example suggests identity or some equivalence for whitespace.\textbackslash{}nThe second example suggests that 'a{-}b' and 'n{-}o\_alt' are not equivalent.\textbackslash{}nThe third example is very interesting: 'a' and 'n' are different characters, yet it returns True. This strongly implies the function is NOT checking for string equality.\textbackslash{}n\textbackslash{}nLet's look at the characters. 'a' and 'n' are both letters. 'a{-}b' contains letters and a hyphen. 'n{-}o\_alt' contains letters, a hyphen, and an underscore.\textbackslash{}n\textbackslash{}nMaybe it checks if both strings consist only of alphanumeric characters? 'a' is alnum, 'n' is alnum. 'a{-}b' has a hyphen (not alnum in strict sense, but often considered punctuation). 'n{-}o\_alt' has underscore.\textbackslash{}n\textbackslash{}nOr maybe it checks if the strings are 'valid identifiers' or something similar? 'a' is valid. 'n' is valid. 'a{-}b' is not a valid Python identifier because of the hyphen. 'n{-}o\_alt' is not valid because of the hyphen.\textbackslash{}n\textbackslash{}nLet's test this hypothesis: Does} \\
            }
        }
    }
    & & \\ \\

    \theutterance \stepcounter{utterance}  
    & & & \multicolumn{4}{p{0.6\linewidth}}{
        \cellcolor[rgb]{0.9,0.9,0.9}{
            \makecell[{{p{\linewidth}}}]{
                \texttt{\tiny{[P1$\langle$GM]}}
                \texttt{Game Violation: ParseError: Response is not valid JSON.} \\
\texttt{Output ONLY:} \\
\texttt{One JSON with confidence\_score, current\_mode, input\_rationale, and hypothesis keys ...} \\
\texttt{OR} \\
\texttt{SOLVE: ```python ... ```} \\
            }
        }
    }
     \\ \\

    \theutterance \stepcounter{utterance}  
    & \multicolumn{4}{p{0.6\linewidth}}{
        \cellcolor[rgb]{0.561,0.737,0.561}{
            \makecell[{{p{\linewidth}}}]{
                \texttt{\tiny{[P1$\rangle$GM]}}
                \texttt{\{} \\
\texttt{"current\_mode": "probing",} \\
\texttt{"input\_rationale": "I have three data points: ('   ', '   ') {-}> True, ('a{-}b', 'n{-}o\_alt') {-}> False, ('a', 'n') {-}> True. The fact that 'a' and 'n' return True suggests the function is not checking for equality. It might be checking if both strings are valid Python identifiers, or if they are both non{-}empty, or something about their character composition. 'a' and 'n' are valid identifiers. '   ' is not a valid identifier (starts with space). So the identifier hypothesis is weak. Let's try to distinguish between strings that return True and those that return False. 'a' (True), 'n' (True), '   ' (True). 'a{-}b' (False), 'n{-}o\_alt' (False). Maybe it checks if the string contains only alphanumeric characters? 'a' is alnum. 'n' is alnum. '   ' is not alnum. So that's not it. Maybe it checks if the string is a valid variable name? 'a' is valid. 'n' is valid. '   ' is not. So that's not it. Let's try more examples to narrow it down. I will test a string with numbers, a string with underscores, and a string with special characters.",} \\
\texttt{"hypothesis": "unknown",} \\
\texttt{"confidence\_score":} \\
            }
        }
    }
    & & \\ \\

    \theutterance \stepcounter{utterance}  
    & & & \multicolumn{4}{p{0.6\linewidth}}{
        \cellcolor[rgb]{0.9,0.9,0.9}{
            \makecell[{{p{\linewidth}}}]{
                \texttt{\tiny{[P1$\langle$GM]}}
                \texttt{Game Violation: ParseError: Response is not valid JSON.} \\
\texttt{Output ONLY:} \\
\texttt{One JSON with confidence\_score, current\_mode, input\_rationale, and hypothesis keys ...} \\
\texttt{OR} \\
\texttt{SOLVE: ```python ... ```} \\
            }
        }
    }
     \\ \\

    \theutterance \stepcounter{utterance}  
    & \multicolumn{4}{p{0.6\linewidth}}{
        \cellcolor[rgb]{0.561,0.737,0.561}{
            \makecell[{{p{\linewidth}}}]{
                \texttt{\tiny{[P1$\rangle$GM]}}
                \texttt{\{} \\
\texttt{"current\_mode": "probing",} \\
\texttt{"input\_rationale": "I need more data points to distinguish between potential hypotheses. The current examples are: ('   ', '   ') {-}> True, ('a{-}b', 'n{-}o\_alt') {-}> False, ('a', 'n') {-}> True. I suspect the function might be checking if both strings are valid Python identifiers, or perhaps if they are both 'simple' strings (e.g., no special characters). 'a' and 'n' are valid identifiers. '   ' is not. So the identifier hypothesis is likely incorrect. Let's try a string that is a valid identifier but contains numbers, and a string with an underscore.",} \\
\texttt{"hypothesis": "unknown",} \\
\texttt{"confidence\_score": 10} \\
\texttt{\}} \\
            }
        }
    }
    & & \\ \\

    \theutterance \stepcounter{utterance}  
    & & & \multicolumn{4}{p{0.6\linewidth}}{
        \cellcolor[rgb]{0.9,0.9,0.9}{
            \makecell[{{p{\linewidth}}}]{
                \texttt{\tiny{[P1$\langle$GM]}}
                \texttt{OUTPUT: (('aba', 'non\_alt'), False)} \\
            }
        }
    }
     \\ \\

    \theutterance \stepcounter{utterance}  
    & \multicolumn{4}{p{0.6\linewidth}}{
        \cellcolor[rgb]{0.561,0.737,0.561}{
            \makecell[{{p{\linewidth}}}]{
                \texttt{\tiny{[P1$\rangle$GM]}}
                \texttt{\{} \\
\texttt{"current\_mode": "probing",} \\
\texttt{"input\_rationale": "New data point: ('aba', 'non\_alt') {-}> False. \textbackslash{}nPrevious points:\textbackslash{}n1. ('   ', '   ') {-}> True\textbackslash{}n2. ('a{-}b', 'n{-}o\_alt') {-}> False\textbackslash{}n3. ('a', 'n') {-}> True\textbackslash{}n4. ('aba', 'non\_alt') {-}> False\textbackslash{}n\textbackslash{}nLet's analyze the True cases:\textbackslash{}n{-} '   ' and '   ': Both are whitespace only.\textbackslash{}n{-} 'a' and 'n': Both are single lowercase letters.\textbackslash{}n\textbackslash{}nLet's analyze the False cases:\textbackslash{}n{-} 'a{-}b' and 'n{-}o\_alt': 'a{-}b' has a hyphen. 'n{-}o\_alt' has hyphen and underscore.\textbackslash{}n{-} 'aba' and 'non\_alt': 'aba' is all letters. 'non\_alt' is all letters but has an underscore? No, 'non\_alt' has an underscore. 'aba' does not.\textbackslash{}n\textbackslash{}nWait, 'aba' is all alphabetic. 'non\_alt' contains an underscore. \textbackslash{}n'a' is alphabetic. 'n' is alphabetic.\textbackslash{}n'   ' is whitespace.\textbackslash{}n'a{-}b' contains a hyphen.\textbackslash{}n\textbackslash{}nHypothesis: The function returns True if both strings consist entirely of alphabetic characters (isalpha()) OR both consist entirely of whitespace (isspace())? \textbackslash{}nLet's check:\textbackslash{}n{-}} \\
            }
        }
    }
    & & \\ \\

    \theutterance \stepcounter{utterance}  
    & & & \multicolumn{4}{p{0.6\linewidth}}{
        \cellcolor[rgb]{0.9,0.9,0.9}{
            \makecell[{{p{\linewidth}}}]{
                \texttt{\tiny{[P1$\langle$GM]}}
                \texttt{Game Violation: ParseError: Response is not valid JSON.} \\
\texttt{Output ONLY:} \\
\texttt{One JSON with confidence\_score, current\_mode, input\_rationale, and hypothesis keys ...} \\
\texttt{OR} \\
\texttt{SOLVE: ```python ... ```} \\
            }
        }
    }
     \\ \\

    \theutterance \stepcounter{utterance}  
    & \multicolumn{4}{p{0.6\linewidth}}{
        \cellcolor[rgb]{0.561,0.737,0.561}{
            \makecell[{{p{\linewidth}}}]{
                \texttt{\tiny{[P1$\rangle$GM]}}
                \texttt{\{} \\
\texttt{"current\_mode": "probing",} \\
\texttt{"input\_rationale": "I have four data points: ('   ', '   ') {-}> True, ('a{-}b', 'n{-}o\_alt') {-}> False, ('a', 'n') {-}> True, ('aba', 'non\_alt') {-}> False. \textbackslash{}nTrue cases: '   ' (whitespace), 'a' (alpha), 'n' (alpha).\textbackslash{}nFalse cases: 'a{-}b' (has hyphen), 'n{-}o\_alt' (has hyphen/underscore), 'aba' (alpha), 'non\_alt' (has underscore).\textbackslash{}n\textbackslash{}nWait, 'aba' is all alpha, yet it returns False when paired with 'non\_alt'. 'non\_alt' has an underscore. 'aba' does not. \textbackslash{}n'a' and 'n' are both alpha and return True.\textbackslash{}n'   ' and '   ' are both whitespace and return True.\textbackslash{}n'a{-}b' has a hyphen. 'n{-}o\_alt' has hyphen and underscore. Returns False.\textbackslash{}n\textbackslash{}nMaybe the function checks if both strings are 'simple' in some way? Or maybe it checks if they are both valid Python identifiers? 'a' is valid. 'n' is valid. '   ' is not valid. So that's out.\textbackslash{}n\textbackslash{}nLet's look at 'aba' vs 'non\_alt'. 'aba' is a palindrome? 'non\_alt' is not. 'a' is} \\
            }
        }
    }
    & & \\ \\

    \theutterance \stepcounter{utterance}  
    & & & \multicolumn{4}{p{0.6\linewidth}}{
        \cellcolor[rgb]{0.9,0.9,0.9}{
            \makecell[{{p{\linewidth}}}]{
                \texttt{\tiny{[P1$\langle$GM]}}
                \texttt{Game Violation: ParseError: Response is not valid JSON.} \\
\texttt{Output ONLY:} \\
\texttt{One JSON with confidence\_score, current\_mode, input\_rationale, and hypothesis keys ...} \\
\texttt{OR} \\
\texttt{SOLVE: ```python ... ```} \\
            }
        }
    }
     \\ \\

    \theutterance \stepcounter{utterance}  
    & \multicolumn{4}{p{0.6\linewidth}}{
        \cellcolor[rgb]{0.561,0.737,0.561}{
            \makecell[{{p{\linewidth}}}]{
                \texttt{\tiny{[P1$\rangle$GM]}}
                \texttt{\{} \\
\texttt{"current\_mode": "probing",} \\
\texttt{"input\_rationale": "I have four data points: ('   ', '   ') {-}> True, ('a{-}b', 'n{-}o\_alt') {-}> False, ('a', 'n') {-}> True, ('aba', 'non\_alt') {-}> False. \textbackslash{}nTrue cases: '   ' (whitespace), 'a' (alpha), 'n' (alpha).\textbackslash{}nFalse cases: 'a{-}b' (has hyphen), 'n{-}o\_alt' (has hyphen/underscore), 'aba' (alpha), 'non\_alt' (has underscore).\textbackslash{}n\textbackslash{}nLet's re{-}evaluate 'aba' vs 'non\_alt'. 'aba' is all alpha. 'non\_alt' has an underscore. \textbackslash{}n'a' and 'n' are both alpha. \textbackslash{}n'   ' and '   ' are both whitespace.\textbackslash{}n'a{-}b' has a hyphen. 'n{-}o\_alt' has hyphen and underscore.\textbackslash{}n\textbackslash{}nMaybe the function checks if both strings are 'valid Python identifiers'? \textbackslash{}n'a' is valid. 'n' is valid. {-}> True.\textbackslash{}n'   ' is NOT valid. {-}> But it returned True. So this hypothesis is wrong.\textbackslash{}n\textbackslash{}nMaybe the function checks if both strings are 'alphanumeric'? \textbackslash{}n'a' is alnum. 'n' is alnum. {-}> True.\textbackslash{}n'   ' is NOT alnum. {-}> But it returned True. So} \\
            }
        }
    }
    & & \\ \\

    \theutterance \stepcounter{utterance}  
    & & & \multicolumn{4}{p{0.6\linewidth}}{
        \cellcolor[rgb]{0.9,0.9,0.9}{
            \makecell[{{p{\linewidth}}}]{
                \texttt{\tiny{[P1$\langle$GM]}}
                \texttt{Game Violation: ParseError: Response is not valid JSON.} \\
\texttt{Output ONLY:} \\
\texttt{One JSON with confidence\_score, current\_mode, input\_rationale, and hypothesis keys ...} \\
\texttt{OR} \\
\texttt{SOLVE: ```python ... ```} \\
            }
        }
    }
     \\ \\

    \theutterance \stepcounter{utterance}  
    & \multicolumn{4}{p{0.6\linewidth}}{
        \cellcolor[rgb]{0.561,0.737,0.561}{
            \makecell[{{p{\linewidth}}}]{
                \texttt{\tiny{[P1$\rangle$GM]}}
                \texttt{\{} \\
\texttt{"current\_mode": "probing",} \\
\texttt{"input\_rationale": "I have four data points: ('   ', '   ') {-}> True, ('a{-}b', 'n{-}o\_alt') {-}> False, ('a', 'n') {-}> True, ('aba', 'non\_alt') {-}> False. \textbackslash{}nTrue cases: '   ' (whitespace), 'a' (alpha), 'n' (alpha).\textbackslash{}nFalse cases: 'a{-}b' (has hyphen), 'n{-}o\_alt' (has hyphen/underscore), 'aba' (alpha), 'non\_alt' (has underscore).\textbackslash{}n\textbackslash{}nLet's look at the False cases again. 'aba' is all alpha. 'non\_alt' has an underscore. \textbackslash{}n'a' and 'n' are both alpha. \textbackslash{}n'   ' and '   ' are both whitespace.\textbackslash{}n'a{-}b' has a hyphen. 'n{-}o\_alt' has hyphen and underscore.\textbackslash{}n\textbackslash{}nMaybe the function checks if both strings are 'valid Python identifiers'? \textbackslash{}n'a' is valid. 'n' is valid. {-}> True.\textbackslash{}n'   ' is NOT valid. {-}> But it returned True. So this hypothesis is wrong.\textbackslash{}n\textbackslash{}nMaybe the function checks if both strings are 'alphanumeric'? \textbackslash{}n'a' is alnum. 'n' is alnum. {-}> True.\textbackslash{}n'   ' is NOT alnum. {-}> But it returned True. So this hypothesis is wrong} \\
            }
        }
    }
    & & \\ \\

    \theutterance \stepcounter{utterance}  
    & & & \multicolumn{4}{p{0.6\linewidth}}{
        \cellcolor[rgb]{0.9,0.9,0.9}{
            \makecell[{{p{\linewidth}}}]{
                \texttt{\tiny{[P1$\langle$GM]}}
                \texttt{Game Violation: ParseError: Response is not valid JSON.} \\
\texttt{Output ONLY:} \\
\texttt{One JSON with confidence\_score, current\_mode, input\_rationale, and hypothesis keys ...} \\
\texttt{OR} \\
\texttt{SOLVE: ```python ... ```} \\
            }
        }
    }
     \\ \\

    \theutterance \stepcounter{utterance}  
    & \multicolumn{4}{p{0.6\linewidth}}{
        \cellcolor[rgb]{0.561,0.737,0.561}{
            \makecell[{{p{\linewidth}}}]{
                \texttt{\tiny{[P1$\rangle$GM]}}
                \texttt{\{} \\
\texttt{"current\_mode": "probing",} \\
\texttt{"input\_rationale": "I have four data points: ('   ', '   ') {-}> True, ('a{-}b', 'n{-}o\_alt') {-}> False, ('a', 'n') {-}> True, ('aba', 'non\_alt') {-}> False. \textbackslash{}nTrue cases: '   ' (whitespace), 'a' (alpha), 'n' (alpha).\textbackslash{}nFalse cases: 'a{-}b' (has hyphen), 'n{-}o\_alt' (has hyphen/underscore), 'aba' (alpha), 'non\_alt' (has underscore).\textbackslash{}n\textbackslash{}nLet's look at the False cases again. 'aba' is all alpha. 'non\_alt' has an underscore. \textbackslash{}n'a' and 'n' are both alpha. \textbackslash{}n'   ' and '   ' are both whitespace.\textbackslash{}n'a{-}b' has a hyphen. 'n{-}o\_alt' has hyphen and underscore.\textbackslash{}n\textbackslash{}nMaybe the function checks if both strings are 'valid Python identifiers'? \textbackslash{}n'a' is valid. 'n' is valid. {-}> True.\textbackslash{}n'   ' is NOT valid. {-}> But it returned True. So this hypothesis is wrong.\textbackslash{}n\textbackslash{}nMaybe the function checks if both strings are 'alphanumeric'? \textbackslash{}n'a' is alnum. 'n' is alnum. {-}> True.\textbackslash{}n'   ' is NOT alnum. {-}> But it returned True. So this hypothesis is wrong} \\
            }
        }
    }
    & & \\ \\

    \theutterance \stepcounter{utterance}  
    & & & \multicolumn{4}{p{0.6\linewidth}}{
        \cellcolor[rgb]{0.9,0.9,0.9}{
            \makecell[{{p{\linewidth}}}]{
                \texttt{\tiny{[P1$\langle$GM]}}
                \texttt{Game Violation: ParseError: Response is not valid JSON.} \\
\texttt{Output ONLY:} \\
\texttt{One JSON with confidence\_score, current\_mode, input\_rationale, and hypothesis keys ...} \\
\texttt{OR} \\
\texttt{SOLVE: ```python ... ```} \\
            }
        }
    }
     \\ \\

    \theutterance \stepcounter{utterance}  
    & \multicolumn{4}{p{0.6\linewidth}}{
        \cellcolor[rgb]{0.561,0.737,0.561}{
            \makecell[{{p{\linewidth}}}]{
                \texttt{\tiny{[P1$\rangle$GM]}}
                \texttt{\{} \\
\texttt{"current\_mode": "probing",} \\
\texttt{"input\_rationale": "I have four data points: ('   ', '   ') {-}> True, ('a{-}b', 'n{-}o\_alt') {-}> False, ('a', 'n') {-}> True, ('aba', 'non\_alt') {-}> False. \textbackslash{}nTrue cases: '   ' (whitespace), 'a' (alpha), 'n' (alpha).\textbackslash{}nFalse cases: 'a{-}b' (has hyphen), 'n{-}o\_alt' (has hyphen/underscore), 'aba' (alpha), 'non\_alt' (has underscore).\textbackslash{}n\textbackslash{}nLet's look at the False cases again. 'aba' is all alpha. 'non\_alt' has an underscore. \textbackslash{}n'a' and 'n' are both alpha. \textbackslash{}n'   ' and '   ' are both whitespace.\textbackslash{}n'a{-}b' has a hyphen. 'n{-}o\_alt' has hyphen and underscore.\textbackslash{}n\textbackslash{}nMaybe the function checks if both strings are 'valid Python identifiers'? \textbackslash{}n'a' is valid. 'n' is valid. {-}> True.\textbackslash{}n'   ' is NOT valid. {-}> But it returned True. So this hypothesis is wrong.\textbackslash{}n\textbackslash{}nMaybe the function checks if both strings are 'alphanumeric'? \textbackslash{}n'a' is alnum. 'n' is alnum. {-}> True.\textbackslash{}n'   ' is NOT alnum. {-}> But it returned True. So this hypothesis is wrong} \\
            }
        }
    }
    & & \\ \\

    \theutterance \stepcounter{utterance}  
    & & & \multicolumn{4}{p{0.6\linewidth}}{
        \cellcolor[rgb]{0.9,0.9,0.9}{
            \makecell[{{p{\linewidth}}}]{
                \texttt{\tiny{[P1$\langle$GM]}}
                \texttt{Game Violation: ParseError: Response is not valid JSON.} \\
\texttt{Output ONLY:} \\
\texttt{One JSON with confidence\_score, current\_mode, input\_rationale, and hypothesis keys ...} \\
\texttt{OR} \\
\texttt{SOLVE: ```python ... ```} \\
            }
        }
    }
     \\ \\

    \theutterance \stepcounter{utterance}  
    & \multicolumn{4}{p{0.6\linewidth}}{
        \cellcolor[rgb]{0.561,0.737,0.561}{
            \makecell[{{p{\linewidth}}}]{
                \texttt{\tiny{[P1$\rangle$GM]}}
                \texttt{\{} \\
\texttt{"current\_mode": "probing",} \\
\texttt{"input\_rationale": "I have four data points: ('   ', '   ') {-}> True, ('a{-}b', 'n{-}o\_alt') {-}> False, ('a', 'n') {-}> True, ('aba', 'non\_alt') {-}> False. \textbackslash{}nTrue cases: '   ' (whitespace), 'a' (alpha), 'n' (alpha).\textbackslash{}nFalse cases: 'a{-}b' (has hyphen), 'n{-}o\_alt' (has hyphen/underscore), 'aba' (alpha), 'non\_alt' (has underscore).\textbackslash{}n\textbackslash{}nLet's look at the False cases again. 'aba' is all alpha. 'non\_alt' has an underscore. \textbackslash{}n'a' and 'n' are both alpha. \textbackslash{}n'   ' and '   ' are both whitespace.\textbackslash{}n'a{-}b' has a hyphen. 'n{-}o\_alt' has hyphen and underscore.\textbackslash{}n\textbackslash{}nMaybe the function checks if both strings are 'valid Python identifiers'? \textbackslash{}n'a' is valid. 'n' is valid. {-}> True.\textbackslash{}n'   ' is NOT valid. {-}> But it returned True. So this hypothesis is wrong.\textbackslash{}n\textbackslash{}nMaybe the function checks if both strings are 'alphanumeric'? \textbackslash{}n'a' is alnum. 'n' is alnum. {-}> True.\textbackslash{}n'   ' is NOT alnum. {-}> But it returned True. So this hypothesis is wrong} \\
            }
        }
    }
    & & \\ \\

    \theutterance \stepcounter{utterance}  
    & & & \multicolumn{4}{p{0.6\linewidth}}{
        \cellcolor[rgb]{0.9,0.9,0.9}{
            \makecell[{{p{\linewidth}}}]{
                \texttt{\tiny{[P1$\langle$GM]}}
                \texttt{Game Violation: ParseError: Response is not valid JSON.} \\
\texttt{Output ONLY:} \\
\texttt{One JSON with confidence\_score, current\_mode, input\_rationale, and hypothesis keys ...} \\
\texttt{OR} \\
\texttt{SOLVE: ```python ... ```} \\
            }
        }
    }
     \\ \\

    \theutterance \stepcounter{utterance}  
    & \multicolumn{4}{p{0.6\linewidth}}{
        \cellcolor[rgb]{0.561,0.737,0.561}{
            \makecell[{{p{\linewidth}}}]{
                \texttt{\tiny{[P1$\rangle$GM]}}
                \texttt{\{} \\
\texttt{"current\_mode": "probing",} \\
\texttt{"input\_rationale": "I have four data points: ('   ', '   ') {-}> True, ('a{-}b', 'n{-}o\_alt') {-}> False, ('a', 'n') {-}> True, ('aba', 'non\_alt') {-}> False. \textbackslash{}nTrue cases: '   ' (whitespace), 'a' (alpha), 'n' (alpha).\textbackslash{}nFalse cases: 'a{-}b' (has hyphen), 'n{-}o\_alt' (has hyphen/underscore), 'aba' (alpha), 'non\_alt' (has underscore).\textbackslash{}n\textbackslash{}nLet's look at the False cases again. 'aba' is all alpha. 'non\_alt' has an underscore. \textbackslash{}n'a' and 'n' are both alpha. \textbackslash{}n'   ' and '   ' are both whitespace.\textbackslash{}n'a{-}b' has a hyphen. 'n{-}o\_alt' has hyphen and underscore.\textbackslash{}n\textbackslash{}nMaybe the function checks if both strings are 'valid Python identifiers'? \textbackslash{}n'a' is valid. 'n' is valid. {-}> True.\textbackslash{}n'   ' is NOT valid. {-}> But it returned True. So this hypothesis is wrong.\textbackslash{}n\textbackslash{}nMaybe the function checks if both strings are 'alphanumeric'? \textbackslash{}n'a' is alnum. 'n' is alnum. {-}> True.\textbackslash{}n'   ' is NOT alnum. {-}> But it returned True. So this hypothesis is wrong} \\
            }
        }
    }
    & & \\ \\

    \theutterance \stepcounter{utterance}  
    & & & \multicolumn{4}{p{0.6\linewidth}}{
        \cellcolor[rgb]{0.9,0.9,0.9}{
            \makecell[{{p{\linewidth}}}]{
                \texttt{\tiny{[P1$\langle$GM]}}
                \texttt{Game Violation: ParseError: Response is not valid JSON.} \\
\texttt{Output ONLY:} \\
\texttt{One JSON with confidence\_score, current\_mode, input\_rationale, and hypothesis keys ...} \\
\texttt{OR} \\
\texttt{SOLVE: ```python ... ```} \\
            }
        }
    }
     \\ \\

    \theutterance \stepcounter{utterance}  
    & \multicolumn{4}{p{0.6\linewidth}}{
        \cellcolor[rgb]{0.561,0.737,0.561}{
            \makecell[{{p{\linewidth}}}]{
                \texttt{\tiny{[P1$\rangle$GM]}}
                \texttt{\{} \\
\texttt{"current\_mode": "probing",} \\
\texttt{"input\_rationale": "I have four data points: ('   ', '   ') {-}> True, ('a{-}b', 'n{-}o\_alt') {-}> False, ('a', 'n') {-}> True, ('aba', 'non\_alt') {-}> False. \textbackslash{}nTrue cases: '   ' (whitespace), 'a' (alpha), 'n' (alpha).\textbackslash{}nFalse cases: 'a{-}b' (has hyphen), 'n{-}o\_alt' (has hyphen/underscore), 'aba' (alpha), 'non\_alt' (has underscore).\textbackslash{}n\textbackslash{}nLet's look at the False cases again. 'aba' is all alpha. 'non\_alt' has an underscore. \textbackslash{}n'a' and 'n' are both alpha. \textbackslash{}n'   ' and '   ' are both whitespace.\textbackslash{}n'a{-}b' has a hyphen. 'n{-}o\_alt' has hyphen and underscore.\textbackslash{}n\textbackslash{}nMaybe the function checks if both strings are 'valid Python identifiers'? \textbackslash{}n'a' is valid. 'n' is valid. {-}> True.\textbackslash{}n'   ' is NOT valid. {-}> But it returned True. So this hypothesis is wrong.\textbackslash{}n\textbackslash{}nMaybe the function checks if both strings are 'alphanumeric'? \textbackslash{}n'a' is alnum. 'n' is alnum. {-}> True.\textbackslash{}n'   ' is NOT alnum. {-}> But it returned True. So this hypothesis is wrong} \\
            }
        }
    }
    & & \\ \\

    \theutterance \stepcounter{utterance}  
    & & & \multicolumn{2}{p{0.3\linewidth}}{
        \cellcolor[rgb]{0.9,0.9,0.9}{
            \makecell[{{p{\linewidth}}}]{
                \texttt{\tiny{[GM$|$GM]}}
                \texttt{Game Verdict: LOSS (turn limit reached)} \\
            }
        }
    }
    & & \\ \\

\end{supertabular}
}

\end{document}
